\begin{definition}[Hasse functions]\label{mod:finitefield-hassefunction}
Let $k$ be a field and let
$\psi\in\operatorname{FiniteAddChar}(k)$.  A Hasse function for $\psi$
is a function $\phi:k\to\mathbf C$ such that $\phi(x)\ne0$ for every
$x$ and
\[
 \phi(x+y)=\phi(x)\phi(y)\psi(xy).
\]
This is the Lean structure $\operatorname{HasseFunction}(k,\psi)$; in
particular, the polar phase has the manuscript's positive sign.
The defining law and nonvanishing imply $\phi(0)=1$.

For $a\in k$, define the translated refinement
\[
 (T_a\phi)(x)=\phi(x)\psi(ax).
\]
It is again a Hasse function for $\psi$, and Lean proves the exact
representative-change identity
\[
 (T_a\phi)(x)=\phi(a)^{-1}\phi(x+a).
\]
For every $n\in\mathbf N$, iteration of the refinement law gives the
scaling formula used in the Hasse-lift calculation:
\[
 \phi(n x)=\phi(x)^n
   \psi\!\left(\binom n2x^2\right).
\]
For $c\in k$, the rescaled function $x\mapsto\phi(cx)$ is a Hasse
function for $\psi^{c^2}(z)=\psi(c^2z)$.

Suppose now that $k$ is finite.  Every value of $\phi$ has absolute
value one.  With
\[
 S(\phi)=\sum_{x\in k}\phi(x),
 \qquad H(\phi)=\Aphase(S(\phi)),
\]
Lean proves
\[
 S(T_a\phi)=\phi(a)^{-1}S(\phi),
 \qquad
 S(x\mapsto\phi(cx))=S(\phi)\quad(c\ne0).
\]
If $\psi\ne1$, additive orthogonality gives the exact identities
\[
 \overline{S(\phi)}S(\phi)=|k|,
 \qquad |S(\phi)|^2=|k|,
 \qquad |S(\phi)|=\sqrt{|k|}.
\]
Thus $S(\phi)\ne0$, and translation transports its phase with no
discarded factor:
\[
 H(T_a\phi)=\phi(a)^{-1}H(\phi).
\]

Finally, if $\operatorname{char}k\ne2$, Lean constructs the standard
quadratic Hasse function
\[
 x\longmapsto\psi\!\left(\frac a2x^2+bx\right)
\]
for the polar character $\psi^a$.  Its sum and its total phase are
definitionally the previously certified
$S_\psi(a,b)$ and $Q_\psi(a,b)$.
\LeanFile{LanglandsFirstMainLemma/FiniteField/HasseFunction.lean}
\lean{LanglandsFirstMainLemma.HasseFunction}
\lean{LanglandsFirstMainLemma.HasseFunction.map_zero}
\lean{LanglandsFirstMainLemma.HasseFunction.translate}
\lean{LanglandsFirstMainLemma.HasseFunction.translate_apply_eq_inv_mul}
\lean{LanglandsFirstMainLemma.HasseFunction.map_nsmul}
\lean{LanglandsFirstMainLemma.HasseFunction.scale}
\lean{LanglandsFirstMainLemma.HasseFunction.norm_apply}
\lean{LanglandsFirstMainLemma.HasseFunction.sum_translate}
\lean{LanglandsFirstMainLemma.HasseFunction.sum_scale}
\lean{LanglandsFirstMainLemma.HasseFunction.sum_star_mul_self}
\lean{LanglandsFirstMainLemma.HasseFunction.sum_norm_sq}
\lean{LanglandsFirstMainLemma.HasseFunction.sum_ne_zero}
\lean{LanglandsFirstMainLemma.HasseFunction.sum_norm}
\lean{LanglandsFirstMainLemma.HasseFunction.sumPhase_translate}
\lean{LanglandsFirstMainLemma.HasseFunction.quadratic}
\lean{LanglandsFirstMainLemma.HasseFunction.sum_quadratic}
\lean{LanglandsFirstMainLemma.HasseFunction.sumPhase_quadratic}
\leanok
\SourceText{Theorem thm:Lamprecht and Theorem thm:Hasse-lift.}
\uses{mod:finitefield-characterapi,mod:finitefield-quadraticphase}
\FormalizationContract
\end{definition}
\begin{proof}
The value at zero follows by cancellation from the refinement law.
The translation identity is that same law with one argument fixed, and
the natural-number formula follows by induction using Pascal's identity.
Argument scaling changes $xy$ to $c^2xy$, while a nonzero $c$ merely
reindexes the finite sum.

To prove unimodularity, specialize the natural-number formula to the
characteristic of $k$ and take absolute values.  Expanding
$\overline{S(\phi)}S(\phi)$, reindexing by $x=y+h$, and applying the
refinement law leaves
\[
 \sum_{h\in k}\phi(h)\sum_{y\in k}\psi(hy).
\]
Nontrivial additive-character orthogonality kills every term except
$h=0$, proving the magnitude and nonvanishing statements.  The phase
identity then transports the norm-one multiplier $\phi(a)^{-1}$ through
the total phase.  In odd characteristic, direct expansion of
$(x+y)^2$ verifies the stated quadratic example and its polar character.
\end{proof}
